\subsection{Différence d'héritage net}
\label{sec:diff_heritage}

Pour évaluer l'avantage relatif de chaque enveloppe, il ne suffit pas de comparer les taux de succession. Une analyse complète doit prendre en compte quatre facteurs clés qui déterminent l'héritage net final pour chaque support.
\begin{itemize}
    \item \textbf{Les frais de gestion :} L'assurance-vie supporte des frais de gestion annuels sur l'encours (typiquement entre 0,5\% et 1\%), qui réduisent mécaniquement le capital transmis. Le CTO, en l'absence de frais d'enveloppe, permet une capitalisation plus performante à rendement brut identique.
    
    \item \textbf{Les prélèvements sociaux (PS) :} Au décès, les gains non encore taxés d'une assurance-vie sont soumis aux prélèvements sociaux (17,2\%). Ce prélèvement, qui n'existe pas pour le CTO, ampute une part significative de la performance.
    
    \item \textbf{La purge des plus-values latentes :} Le CTO bénéficie d'un avantage fiscal majeur : au décès, les plus-values latentes sont "purgées". Les héritiers reçoivent les titres valorisés au jour du décès, et l'impôt sur le revenu qui aurait été dû en cas de vente par le défunt est effacé. Cet avantage est d'autant plus important que la durée de détention et le rendement ont été élevés.
    
    \item \textbf{Le barème successoral :} C'est le seul domaine où l'assurance-vie (pour les versements avant 70 ans) dispose d'un avantage théorique, avec un abattement de 152 500 € par bénéficiaire et des taux spécifiques (20\% et 31,25\%). Le CTO est quant à lui intégré à la succession classique, soumis au barème général après l'abattement de 100 000 € en ligne directe.
\end{itemize}

Trop souvent, le débat public réduit la comparaison à une simple opposition des barèmes successoraux. Cette vision est trompeuse. Comme le montreront les simulations, l'avantage fiscal du CTO (absence de frais d'enveloppe et purge des plus-values) l'emporte très souvent sur l'avantage successoral de l'assurance-vie, dont l'efficacité est diminuée par les frais et les prélèvements sociaux. La question n'est donc pas de savoir si l'assurance-vie est taxée, mais si son régime spécifique reste avantageux une fois tous les prélèvements pris en compte.
